% 01_introduction.tex — Rail Paper: Introduction
% ============================================================

\section{Introduction}
\label{sec:intro}

% ------------------------------------------------------------------
%  The 14.5 dB paradox
% ------------------------------------------------------------------

In 2022, the Mask-guided Spectral-wise Transformer (MST-L)~\cite{cai2022mst}
achieved 34.81\,dB on the standard CASSI benchmark---a number that
represents the culmination of nearly a decade of algorithmic progress in
hyperspectral image reconstruction.
Yet when the coded aperture mask is shifted by a single pixel---a
perturbation well within the manufacturing and alignment tolerances of
any real optical system---the same method collapses to 18.09\,dB, a loss
of \textbf{16.72\,dB}.
Under a 1-pixel-shift-plus-rotation model, neural solvers lose
10--17\,dB (\cref{tab:sensitivity}).
To put this in perspective, the entire history of deep-learning-based
CASSI reconstruction---from GAP-TV's 24.22\,dB in 2016 to MST-L's
34.81\,dB in 2022---spans only 10.59\,dB.
A single pixel of operator mismatch erases more than a decade of
algorithmic improvement.

This is not an isolated failure.
Across modalities---coded aperture snapshot spectral imaging (CASSI),
coded aperture compressive temporal imaging (CACTI), single-pixel
cameras (SPC)---the pattern repeats: the most powerful neural solvers
exhibit the \emph{largest} sensitivity to forward-model error.
Classical methods such as GAP-TV lose only 4.56\,dB under the same
mismatch, precisely because they lack the capacity to overfit the
operator.
The paradox is stark: the field has been building faster and faster
trains while ignoring the fact that the rails are broken.

% ------------------------------------------------------------------
%  The thesis: wrong layer
% ------------------------------------------------------------------

\paragraph{The thesis.}
We argue that the computational imaging community is optimizing the
wrong layer.
The bottleneck is not the \emph{solver}---the reconstruction algorithm
that maps measurements to images---but the \emph{infrastructure} around
it: the evaluation protocols, the physics representations, the
calibration pipelines, and the benchmarks that determine what counts as
progress.
In the language of industrial revolutions, solvers are
\emph{trains}---visible, glamorous, and destined to commoditize.
The compounding, durable value lies in the \emph{rails}: the standards,
metrics, and institutional infrastructure that make it possible for any
train to run reliably.

This paper presents the \textbf{\pwm{} (PWM)} as precisely such a rail
for computational imaging.
PWM is not a new reconstruction algorithm.
It is an \emph{evaluation harness}---a standardized infrastructure layer
that makes it possible to measure, compare, and improve reconstruction
methods under realistic physical conditions.

% ------------------------------------------------------------------
%  SolveEverything.org framing
% ------------------------------------------------------------------

\paragraph{SolveEverything framing.}
Our design follows the SolveEverything.org framework~\cite{wissner2024solveeverything},
which identifies a recurring pattern across domains---from protein
folding to chip design to weather prediction---in which transformative
progress requires not just better models but better infrastructure.
SolveEverything characterizes the pre-infrastructure phase as
\emph{The Muddle} (maturity levels L0--L1): a regime in which
practitioners cannot agree on what to measure, how to measure it, or
what the measurements mean.

Computational imaging is firmly in The Muddle.
Every laboratory uses different test scenes, different noise models,
different mismatch assumptions (if any), and different metrics.
A reconstruction method that reports 34\,dB on one benchmark may achieve
22\,dB on another, and neither number tells us anything about real-world
deployment.
There is no standard forward-model representation, no common evaluation
protocol, and no mechanism for prospective (rather than retrospective)
assessment.
The consequence is a literature of mutually incomparable results, an
illusion of progress driven by leaderboard optimization, and a
sim-to-real gap that remains invisible until hardware deployment.

% ------------------------------------------------------------------
%  Paper overview
% ------------------------------------------------------------------

\paragraph{Paper overview.}
We present PWM as the evaluation harness that moves computational
imaging from The Muddle toward industrialization.
The system comprises four interlocking components (\cref{fig:stack}):

\begin{enumerate}
  \item \textbf{OperatorGraph IR}---a universal directed acyclic graph
        (DAG) representation for forward measurement operators that
        spans 64 modalities across 5 physical carriers (photons,
        electrons, spins, acoustic waves, particles) using 89 validated
        templates.
  \item \textbf{4-Scenario Evaluation Protocol}---a standardized
        framework that measures every solver under Ideal, Assumed
        (mismatched), Corrected (calibrated), and Oracle conditions,
        isolating the contribution of operator knowledge from solver
        architecture.
  \item \textbf{LIP-Arena}---a prospective evaluation competition using
        a Commit-Measure-Score protocol in which test data is generated
        \emph{after} the submission deadline, eliminating memorization
        and overfitting to known benchmarks.
  \item \textbf{Red Team module}---an adversarial evaluation layer that
        probes submitted methods against novel mismatches, compound
        perturbations, out-of-family physics, distribution shifts,
        compute traps, and gate-flip scenarios.
\end{enumerate}

% ------------------------------------------------------------------
%  Contributions
% ------------------------------------------------------------------

\paragraph{Contributions.}
Our specific contributions are:

\begin{enumerate}
  \item \textbf{SolveEverything mapping to imaging.}
        We provide the first systematic mapping of the SolveEverything
        Industrial Intelligence Stack---all 9 layers and the L0--L5
        maturity ladder---to the computational imaging domain, producing
        a concrete roadmap from the current L1 state to full
        industrialization (\cref{sec:framework}).

  \item \textbf{OperatorGraph IR as universal physics representation.}
        We define a graph-based intermediate representation for forward
        measurement operators that unifies 64 modalities under a single
        formalism, with 89 validated templates supporting automatic
        differentiation, adjoint computation, and mismatch injection
        (\cref{sec:architecture}).

  \item \textbf{26-modality benchmark evidence.}
        We present the first cross-modality sensitivity analysis
        demonstrating that neural solvers systematically amplify
        operator mismatch, with degradation ranging from 4.56\,dB
        (GAP-TV) to 16.72\,dB (MST-L) on CASSI alone
        (\cref{sec:problem}).

  \item \textbf{Operator correction across 9 validated modalities
        (16 registered).}
        We show that calibrating the forward operator---even with
        classical methods---recovers a substantial fraction of mismatch
        loss across 9 numerically validated modalities (with 7 additional
        modalities registered for future evaluation), confirming that the
        dominant failure mode in deployed imaging systems is not solver
        quality but operator fidelity.

  \item \textbf{LIP-Arena prospective evaluation protocol.}
        We introduce the first Commit-Measure-Score competition for
        computational imaging, with anti-Goodhart scoring, safety
        brakes, and outcome-based governance, providing a sustainable
        institutional mechanism for measuring genuine progress.
\end{enumerate}

The remainder of this paper is organized as follows.
\Cref{sec:problem} quantifies the solver-only optimization trap using
CASSI as a case study.
\Cref{sec:framework} maps the SolveEverything framework to
computational imaging.
\Cref{sec:architecture} presents the PWM architecture in detail.
Subsequent sections describe the Triad Law diagnostic framework,
multi-agent orchestration, experiments, and implications for the field.
